% Solutions to the Chapter 2 exercises that are not code, from lectures/L02/appendix/c_exercises.md.
% Exercises 2.6 and 2.7 are Code and are not answered; see chapter.tex for why.

\section{Chapter 2: The Command Line, and the System Underneath It}
\label{sol:sec:c2}
Answers to the exercises in \secref{c2:sec:exercises}. Exercises~\ref{c2:ex:6} and~\ref{c2:ex:7}
are Code, and are checked on the target: the first by \code{make test L=L02}, as its Check yourself
line says, and the second by the machine itself.

% ------------------------------------------------------------------------------------------
\solution[fillaccent2]{c2:ex:1}{What the shell does and what it does not}{Recall}

\xpart{a} The shell. Pathname expansion happens before \code{exec}, so the program receives the
list of names in \code{argv} and never sees the \code{*}.

\xpart{b} The program. To the shell \code{-l} is a word like any other; each program decides for
itself what its options mean, which is why \code{-l} means ``long'' to \code{ls} and ``lines'' to
\code{wc}.

\xpart{c} The shell. It opens \file{out.txt} and \code{dup2}s it onto descriptor 1 before
\code{exec}; the program writes to descriptor 1 as it always does and cannot tell.

\xpart{d} The shell. \code{cd} is a builtin, because a child can change only its own directory and
never its parent's (\secref{c2:app:a1}).

\xpart{e} Whichever of them made the call that failed, and for a file named as an argument that is
the program: \code{cat missing} gets \code{ENOENT} from its own \code{open} and prints the message
itself. If the missing file is a redirection, \code{< missing}, the shell made the call and the
shell reports it. The words are the C library's text for \code{ENOENT} either way.

\xpart{f} The shell. Parameter expansion replaces \code{$HOME} with its value before the program
starts, and the program receives the value.

\textbf{Then:} a shell leaves a pattern that matches nothing as it is, so \code{grep} receives three
arguments: \code{grep}, \code{foo}, and the literal \code{*.log}. It tries to open a file called
\code{*.log}, fails, prints \code{grep: *.log: No such file or directory}, and exits with status 2,
which means an error, as distinct from 1, which means no match. The message reads like a missing
file, and it is the shell's expansion rule that made one. (Bash can be told otherwise, with
\code{nullglob}; the default, and BusyBox's \code{ash}, pass the pattern through.)

% ------------------------------------------------------------------------------------------
\solution[fillaccent2]{c2:ex:2}{Pipelines}{Recall}

\xpart{a} Three processes, one for each stage, each forked by the shell and then \code{exec}'d, and
two pipes, one between each pair of neighbours. Plumbing them takes four \code{dup2}s: \code{ps}'s
output, \code{grep}'s input and output, and \code{wc}'s input.

\xpart{b} \code{B} never sees end of file. A pipe's read end returns end of file only when every
copy of the write end is closed, and the shell's copy keeps it open after \code{A} has finished and
exited. So \code{B}'s \code{read} blocks forever, and the pipeline hangs, with the shell waiting for
\code{B}. \code{B} is the process that suffers; \code{A} runs and exits normally.

\xpart{c} The pipe's buffer is bounded, 64 KB, sixteen pages of 4 KB. \code{yes} writes until it is
full and then blocks, so it can never hold more memory than that. \code{head} reads one line,
prints it and exits, and then no process holds the read end open. \code{yes}'s next write raises
\code{SIGPIPE}, whose default action is to terminate it, and it dies with status $128 + 13 = 141$
(\secref{c2:app:a3}). Had it ignored \code{SIGPIPE}, the write would have failed with \code{EPIPE}
instead, and a well-behaved \code{yes} would have exited on that.

\xpart{d} A pipeline's exit status is its last stage's: \code{wc -l} ran perfectly well and printed
0, so \code{$?} is 0 and \code{grep}'s 1 is lost (\secref{c2:app:a3}). Add \code{set -o pipefail},
which makes a pipeline's status the last non-zero status among its stages; BusyBox's \code{ash}
supports it. If the count is what you want, test it instead: \code{grep -c} or \code{grep -q} gives
the answer without a pipeline at all.

% ------------------------------------------------------------------------------------------
\solution[fillaccent3]{c2:ex:3}{Reading \code{/proc/stat}}{Hand calculation}

The interval is $6.40 - 3.90 = 2.50$ s, from the first field of \file{/proc/uptime}.

\xpart{a} $1901 - 1626 = 275$ context switches in 2.50 s, so $275 / 2.50 = 110$ per second.

\xpart{b} $71 - 63 = 8$ processes, so $8 / 2.50 = 3.2$ per second. \code{processes} counts forks
since boot, so this is a fork rate.

\xpart{c} $699 - 252 = 447$ units of \code{USER_HZ}, which is $447 / 100 = 4.47$ s. The unit is
\code{USER_HZ}, 1/100 s, because \file{/proc/stat} reports CPU time in it whatever \code{CONFIG_HZ}
is (\secref{c2:app:b2}); dividing by \code{HZ} instead, 250, would give 1.79 s, 2.5 times too small.
\file{/proc/uptime}'s second field, which the kernel computes from the same idle time, agrees: it
rose by $6.99 - 2.52 = 4.47$ s.

\xpart{d} Two CPUs for 2.50 s is $2 \times 2.50 = 5.00$ s of CPU time, 500 units, and the ten fields
of the \code{cpu} line, five of which are not zero, do rise by exactly that:
\[
(19 + 383 + 699 + 81 + 27) - (16 + 341 + 252 + 75 + 25) = 1209 - 709 = 500.
\]
Of it, 4.47 s was idle, $4.47 / 5.00 = 89.4\%$; the other 53 units are user 3, system 42, irq 6 and
softirq 2. Yes, 4.47 s exceeds the 2.50 s interval, and it is not a mistake: the \code{cpu} line
sums over both CPUs, so it can gather up to 5.00 s in 2.50 s. Per CPU, the idle time is
$4.47 / 2 = 2.235$ s of the 2.50.

\xpart{e} None of them. The CPU times are in \code{USER_HZ}, which is 100 whatever \code{CONFIG_HZ}
is, so 447 units are still 4.47 s; \code{ctxt} and \code{processes} are plain counts; and the
interval comes from \file{/proc/uptime}, in seconds. A kernel ticking at 1000 Hz behaves slightly
differently, so the counts themselves could come out different on another run, but no step of the
arithmetic involves \code{HZ}, and noticing that is the point of the question.

% ------------------------------------------------------------------------------------------
\solution[fillaccent3]{c2:ex:4}{Device numbers}{Hand calculation}

\xpart{a} It tells you that one driver handles both: major 1, which \file{/proc/devices} lists as
\code{mem}. Minors 3 and 5 are two devices of that driver's, which it tells apart by minor. It does
not tell you what either device does, which is decided by the driver's code and not by the numbers;
nor that the two share anything beyond the driver; nor that any driver is registered at all, which
only \file{/proc/devices} or opening the node can tell you.

\xpart{b} The tty core, \file{drivers/tty/tty_io.c}. The excerpt lists major 5 as \code{/dev/tty};
the full \file{/proc/devices} lists it three times, as \code{/dev/tty}, \code{/dev/console} and
\code{/dev/ptmx}, because the tty core registers each of minors 0, 1 and 2 separately. Minor 1 is
the console, and it is an indirection rather than a device of its own: opening it opens whichever
tty the kernel's console is on, here the PL011 UART's \file{ttyAMA0}, named by \code{console=ttyAMA0}
on the command line (\secref{c2:app:b2}).

\xpart{c} The \code{mknod} succeeds: it writes a directory entry holding \code{c}, 204 and 64, and
consults no driver. Opening the node succeeds as well, because major 204 is claimed by the PL011
driver, \code{ttyAMA}, and minor 64 is its first port: \file{/dev/spare} is a second name for
\file{/dev/ttyAMA0}, the console you are typing on. They are not the same question. The first is
about the filesystem, and the second is about whether a driver has registered the number
(\secref{c2:app:b5}); here both answers happen to be yes.

\xpart{d} The \code{mknod} succeeds, for the same reason. Opening it fails with \code{ENXIO}, ``No
such device or address'', because no driver has registered major 42 (\secref{c2:app:b5}). Here the
answers differ, which is what makes them two questions.

\xpart{e} What is gone is one directory entry in devtmpfs: the name. The \code{mem} driver and its
minor 3 are still registered, any process that already had \file{/dev/null} open goes on using it,
and \code{mknod /dev/null c 1 3} with a \code{chmod 666} gives the device its name back. Until then
the loss is worse than it sounds: the next \code{> /dev/null} a root shell runs creates a regular
file of that name, which then collects everything that was meant to be thrown away.

% ------------------------------------------------------------------------------------------
\solution[fillaccent4]{c2:ex:5}{Which file answers the question}{Design}

\xpart{a} \file{/proc/version}, whose \code{SMP PREEMPT} names the model the running kernel was
built with (\secref{c2:app:b2}). The obvious alternative, the \code{.config} on the build machine,
describes whatever you built last, which need not be what is running; \file{/proc/config.gz} is
exact, but exists only on a kernel built with \code{CONFIG_IKCONFIG_PROC}, as this one happens to
be.

\xpart{b} \file{/sys/devices/system/cpu/online}, which holds one value, a list of ranges such as
\code{0-1}, generated from the kernel's own mask of online CPUs. The obvious alternative, counting
the \code{processor} lines in \file{/proc/cpuinfo}, depends on a format written for people and
different on every architecture; it works on arm64, which is why the lab's test uses it, and an
answer that chooses it has to argue that portability does not matter.

\xpart{c} \file{/proc/cmdline}, exactly what the kernel was given. The alternative, the
\code{Kernel command line:} line in \code{dmesg}, lives in a ring buffer that a machine that has been
up long enough has overwritten.

\xpart{d} \file{/proc/interrupts}: the \code{arch_timer} row, in the \code{CPU1} column, which is one
count per CPU per interrupt source. The alternative, the \code{intr} line of \file{/proc/stat},
holds one total per interrupt number summed over every CPU, with no names, so it cannot answer the
question at all.

\xpart{e} \file{/proc/iomem}, read as root: the physical address map as a tree, with each region a
driver has claimed nested under the device's own region with the claimant's name. Read as anyone
else, every address in it is zero. The alternative, the \code{reg} properties under
\file{/sys/firmware/devicetree/base}, says where devices are, not whether any driver has claimed
them: \code{qa-dev}'s node is there from the first boot, and no driver claims it until
Chapter~\ref{c6:ch}.

\xpart{f} \file{/proc/1/comm}, the name the kernel keeps for PID 1. While \file{/init} runs it is
\code{init}, and at a \code{make boot} prompt it is \code{sh}, because the script has
\code{exec}'d \file{/bin/sh} and PID 1 is now the shell. The alternatives are worse:
\file{/proc/1/cmdline} holds the argument vector, which for a script names the interpreter,
\code{/bin/sh /init}, and which a program can rewrite; and \file{/proc/1/exe} points at
\file{/bin/busybox}, as it does for every process on this target.

\xpart{g} \file{/proc/56/fd/}: one symbolic link per open descriptor, named by its number and
pointing at what is open, so \code{ls -l /proc/56/fd} answers it, and \file{/proc/56/fdinfo/}
adds each one's offset and flags. The alternative is \code{lsof}, which is a program that reads
exactly this directory and formats it; BusyBox's is much reduced, and on a target without it the
directory is still there.

% ------------------------------------------------------------------------------------------
\solution[fillaccent4]{c2:ex:8}{What a two-megabyte userland costs you}{Design}

\xpart{a} \code{ls -l /bin/busybox} gives 2,108,608 bytes and \code{busybox --list | wc -l} gives
402 (\secref{c2:app:a8}). The average is $2{,}108{,}608 / 402 \approx 5{,}245$ bytes, about 5 KB per
command. It is misleading three ways. The applets are nothing like equal: \code{ash}, \code{awk} and
\code{vi} are large, and \code{true} and \code{yes} are a few lines. Most of the binary is shared,
BusyBox's own library and the statically linked C library that every applet uses, so no applet costs
5 KB, and removing half of them would shrink the file by much less than half. And it invites a
comparison with coreutils, about 5 MB for about 100 commands, roughly 50 KB each, that compares
different amounts of function, because BusyBox's applets take fewer options.

\xpart{b} \code{busybox --list | grep -x} and the name prints nothing for four of the five.
\code{strace}, \code{ss} and \code{vmstat}: absent, because no BusyBox applet exists. Each is a
separate program, from strace, iproute2 and procps respectively, that could be cross-compiled and
installed. \code{lsof}: present, a BusyBox applet. \code{journalctl}: absent for the other reason:
it belongs to systemd, this system has no systemd, and so there is no journal for it to read.

\xpart{c} Read \file{/proc} directly, which is what \code{lsof} does anyway (\secref{c2:app:b3}).
Find the process's PID with \code{ps}; \code{ls -l /proc/PID/fd} lists one symbolic link per open
descriptor, with what it points at; \file{/proc/PID/fdinfo/N} gives each one's position and flags;
\file{/proc/PID/maps} shows the files it has mapped, which are not descriptors; and \file{cwd} and
\file{exe} are its working directory and its binary.

\xpart{d} A difference in behaviour between the BusyBox applet you ship and the GNU tool you develop
against, which a script relies on without knowing, and the silent kind is the worst. The course's
example is \code{rmmod}: BusyBox's returns success even when the kernel refused to unload the
module, so a test that checks its exit status passes on a laptop and lies on the target, and
Chapter~\ref{c4:ch}'s lab checks \code{lsmod} instead (\secref{c2:app:a8}).

% ------------------------------------------------------------------------------------------
\solution[fillaccent]{c2:ex:9}{A rate measured three ways}{Cross-check}

\xpart{a} The rate is the difference between the two \code{ctxt} readings divided by the difference
between the two first fields of \file{/proc/uptime}. The numbers are yours. Divide by the uptime
difference rather than by the pause on your watch: the uptimes time the reads, and the watch times
you.

\xpart{d}
\begin{itemize}
  \item \textbf{The spread} is yours to compute, and whatever it is, \code{ctxt-per-s} is a
    property of the moment. \code{ctxt} counts every switch on both CPUs, and on an idle target that
    is dominated by whatever happened to run in that second: the kernel's housekeeping, the serial
    console echoing what you type, and the script's own processes. Two runs a second apart measure
    two different seconds.
  \item \textbf{The digits.} The interval comes from \file{/proc/uptime}, which counts in 10 ms
    steps, so a one-second interval is known to about 1\%, and a count of some tens of switches
    moves by one or two percent when one more lands inside it. The run-to-run spread from c) is the real
    uncertainty, and it is far larger than the digit after the point. An honest presentation is a
    whole number, or a range from several runs, stated with the interval and the conditions:
    ``about so many per second, idle, over one second, on QEMU''.
  \item \textbf{Hand against script.} Expect the longer interval to be steadier: it averages bursts
    over more time, and the fixed errors, the 10 ms steps and the switches the reading commands
    cause, are a smaller fraction of it. One reason it may not hold: during a hand-timed pause you
    type the second command, and every keystroke on the serial console is an interrupt and a
    wake-up of the shell, so the hand figure includes you.
  \item \textbf{Without the \code{sleep}.} The rate rises, and stops describing the machine. The
    switches the script itself causes between its two samples, forking \code{cut} and \code{awk},
    their \code{exec}s, their exits and the shell's \code{wait}s, are about the same number
    whatever the interval, so dividing them by 0.14 s instead of about 1 s inflates them roughly
    sevenfold, and they swamp the background. The interval is also only 14 steps of
    \file{/proc/uptime}, so it is uncertain by up to 10 ms in 140, about 7\%. A short interval is
    worse, not better, because a rate is an average over the interval, and there is no ``now'' to
    get closer to: context switches are discrete events, and a rate at an instant does not exist.
\end{itemize}

\xpart{e} \code{vmstat} is not there: the shell says \code{sh: vmstat: not found}, and the status
is 127. That tells you the userland is BusyBox and nothing else: there is no procps, and BusyBox has
no \code{vmstat} applet (\secref{c2:app:a8}). To check the answer without it, use the number
BusyBox does have under another name: \code{nmeter '\%x'} prints the context-switch rate once a
second, computed from the same \code{ctxt} counter, and \code{nmeter} is in this build. Your hand
figure from a) is a second, independent check, the same counter read by a different route; and
running your script on the host, beside the \code{vmstat} there, checks the method rather than the
number.

\xpart{f} Any two of these:
\begin{itemize}
  \item \textbf{It is a property of a moment, not of the machine.} It is one sample of an idle
    target, dominated by whatever ran in that second, including the measurement itself. Under any
    load, or a second later, it is different, and a datasheet figure has to say under what load and
    by what method it was measured.
  \item \textbf{Its precision is invented.} Three significant figures from a one-second sample whose
    interval is known to 10 ms and whose run-to-run spread is far larger than 0.1 per second; one or
    two significant figures at most.
  \item \textbf{It describes an emulation.} The target is QEMU on some host, and the rate depends
    on that host's speed and load, and on this kernel's configuration; it says nothing about an ARM
    board running the same software.
\end{itemize}
